\chapter{Etapa 1 --- Gramática Livre de Contexto e Autômato de Pilha}

\section{Especificação Formal da GLC ($G$)}

O alfabeto de terminais da Etapa 1 é $\Sigma_{\text{AP}} = \{\texttt{p}, \texttt{\{}, \texttt{\}}, \texttt{(}, \texttt{)}, \texttt{\#}, \texttt{i}, \texttt{s}, \texttt{:}, \texttt{;}\}$. A gramática $G = (V, T, P, S)$ é definida por $V = \{S, B, T, R\}$, $T = \Sigma_{\text{AP}}$, símbolo inicial $S$ e 9 produções LL(1) não-ambíguas:

\begin{align*}
P_1:\quad S &\to \texttt{p\{}\,B\,\texttt{\}\#}\,T\,\texttt{:}\,T\,\texttt{;} \\
P_2:\quad B &\to \texttt{\{}\,B\,\texttt{\}}\,B \qquad\qquad
P_3:\quad B \to \texttt{(}\,B\,\texttt{)}\,B \qquad\qquad
P_4:\quad B \to \varepsilon \\
P_5:\quad T &\to \texttt{i}\,R \qquad\qquad\qquad
P_6:\quad T \to \texttt{s}\,R \\
P_7:\quad R &\to \texttt{i}\,R \qquad\qquad\qquad
P_8:\quad R \to \texttt{s}\,R \qquad\qquad\qquad
P_9:\quad R \to \varepsilon
\end{align*}

\section{Derivação Passo a Passo de uma Cadeia Aceita ($\ge 18$ caracteres)}

Abaixo apresenta-se a derivação mais à esquerda completa da cadeia aceita \texttt{p\{(())\{\}\}\#isisi:isisi;} (22 caracteres, caso \texttt{AP-A01}), obtida de forma determinística pelo motor \texttt{src/av2/glc.py}:

\begin{enumerate}
    \item $S \Rightarrow \texttt{p\{}B\texttt{\}\#}T\texttt{:}T\texttt{;}$ \hfill $[P_1]$
    \item $\Rightarrow \texttt{p\{(}B\texttt{)}B\texttt{\}\#}T\texttt{:}T\texttt{;}$ \hfill $[P_3]$
    \item $\Rightarrow \texttt{p\{((}B\texttt{)}B\texttt{)}B\texttt{\}\#}T\texttt{:}T\texttt{;}$ \hfill $[P_3]$
    \item $\Rightarrow \texttt{p\{(())}B\texttt{)}B\texttt{\}\#}T\texttt{:}T\texttt{;}$ \hfill $[P_4]$
    \item $\Rightarrow \texttt{p\{(())}B\texttt{\}\#}T\texttt{:}T\texttt{;}$ \hfill $[P_4]$
    \item $\Rightarrow \texttt{p\{(())\{}B\texttt{\}}B\texttt{\}\#}T\texttt{:}T\texttt{;}$ \hfill $[P_2]$
    \item $\Rightarrow \texttt{p\{(())\{\}}B\texttt{\}\#}T\texttt{:}T\texttt{;}$ \hfill $[P_4]$
    \item $\Rightarrow \texttt{p\{(())\{\}\}\#}T\texttt{:}T\texttt{;}$ \hfill $[P_4]$
    \item $\Rightarrow \texttt{p\{(())\{\}\}\#i}R\texttt{:}T\texttt{;}$ \hfill $[P_5]$
    \item $\Rightarrow \texttt{p\{(())\{\}\}\#is}R\texttt{:}T\texttt{;}$ \hfill $[P_8]$
    \item $\Rightarrow \texttt{p\{(())\{\}\}\#isi}R\texttt{:}T\texttt{;}$ \hfill $[P_7]$
    \item $\Rightarrow \texttt{p\{(())\{\}\}\#isis}R\texttt{:}T\texttt{;}$ \hfill $[P_8]$
    \item $\Rightarrow \texttt{p\{(())\{\}\}\#isisi}R\texttt{:}T\texttt{;}$ \hfill $[P_7]$
    \item $\Rightarrow \texttt{p\{(())\{\}\}\#isisi:}T\texttt{;}$ \hfill $[P_9]$
    \item $\Rightarrow \texttt{p\{(())\{\}\}\#isisi:i}R\texttt{;}$ \hfill $[P_5]$
    \item $\Rightarrow \texttt{p\{(())\{\}\}\#isisi:is}R\texttt{;}$ \hfill $[P_8]$
    \item $\Rightarrow \texttt{p\{(())\{\}\}\#isisi:isi}R\texttt{;}$ \hfill $[P_7]$
    \item $\Rightarrow \texttt{p\{(())\{\}\}\#isisi:isis}R\texttt{;}$ \hfill $[P_8]$
    \item $\Rightarrow \texttt{p\{(())\{\}\}\#isisi:isisi}R\texttt{;}$ \hfill $[P_7]$
    \item $\Rightarrow \texttt{p\{(())\{\}\}\#isisi:isisi;}$ \hfill $[P_9]$
\end{enumerate}

\section{Especificação Formal do Autômato de Pilha ($M_{\text{AP}}$ --- 9 estados)}

Define-se $M_{\text{AP}} = (Q_{\text{AP}}, \Sigma_{\text{AP}}, \Gamma_{\text{AP}}, \delta_{\text{AP}}, q_0, Z, F_{\text{AP}})$ com:
\begin{itemize}
    \item \textbf{Estados (9 estados $\ge 8$):} $Q_{\text{AP}} = \{q_0, q_1, q_2, q_3, q_4, q_5, q_6, q_7, q_8\}$;
    \item \textbf{Alfabeto de entrada:} $\Sigma_{\text{AP}} = \{\texttt{p}, \texttt{\{}, \texttt{\}}, \texttt{(}, \texttt{)}, \texttt{\#}, \texttt{i}, \texttt{s}, \texttt{:}, \texttt{;}\}$;
    \item \textbf{Alfabeto da pilha:} $\Gamma_{\text{AP}} = \{Z, \texttt{\{}, \texttt{(}\}$ (com símbolo inicial $Z$);
    \item \textbf{Configuração inicial:} $(q_0, x, Z)$;
    \item \textbf{Critério de aceitação:} aceitação por \textbf{estado final} ($F_{\text{AP}} = \{q_8\}$) com \textbf{consumo integral da entrada} (e pilha final igual a $Z$).
\end{itemize}

\begin{table}[ht]
\centering
\caption{Função de transição $\delta_{\text{AP}}(q, a, X) = (q', \gamma)$ do DPDA de 9 estados.}
\small
\begin{tabular}{@{}lllll@{}}
\toprule
Estado & Lê & Desempilha & Destino & Empilha (topo à esq.) \\
\midrule
$q_0$ & \texttt{p} & $Z$ & $q_1$ & $Z$ \\
$q_1$ & \texttt{\{} & $Z$ & $q_2$ & \texttt{\{}$Z$ \\
$q_2$ & \texttt{\{} & \texttt{\{} / \texttt{(} & $q_2$ & \texttt{\{\{} / \texttt{\{(} \\
$q_2$ & \texttt{(} & \texttt{\{} / \texttt{(} & $q_2$ & \texttt{(\{} / \texttt{((} \\
$q_2$ & \texttt{\}} & \texttt{\{} & $q_2$ & $\varepsilon$ \\
$q_2$ & \texttt{)} & \texttt{(} & $q_2$ & $\varepsilon$ \\
$q_2$ & $\varepsilon$ & $Z$ & $q_3$ & $Z$ \\
$q_3$ & \texttt{\#} & $Z$ & $q_4$ & $Z$ \\
$q_4$ & \texttt{i} / \texttt{s} & $Z$ & $q_5$ & $Z$ \\
$q_5$ & \texttt{i} / \texttt{s} & $Z$ & $q_5$ & $Z$ \\
$q_5$ & \texttt{:} & $Z$ & $q_6$ & $Z$ \\
$q_6$ & \texttt{i} / \texttt{s} & $Z$ & $q_7$ & $Z$ \\
$q_7$ & \texttt{i} / \texttt{s} & $Z$ & $q_7$ & $Z$ \\
$q_7$ & \texttt{;} & $Z$ & $q_8$ & $Z$ \\
\bottomrule
\end{tabular}
\end{table}

\section{Justificativa da Equivalência $L(G) = L(M_{\text{AP}})$}

Os estados $q_0$ e $q_1$ impõem exatamente o prefixo \texttt{p\{} de $P_1$, empilhando \texttt{\{} sobre $Z$. Enquanto o escopo principal permanece aberto, a pilha em $q_2$ mantém o invariante de armazenar, acima de $Z$, exatamente a sequência de delimitadores \texttt{\{} e \texttt{(} abertos e ainda não fechados, o que corresponde biunivocamente às derivações de $B$ ($P_2, P_3, P_4$) até o fechamento do \texttt{\}} externo de $P_1$. Nesse exato instante o topo retorna a $Z$, habilitando exclusivamente a transição determinística $(q_2, \varepsilon, Z) \to (q_3, Z)$. Por fim, $q_3 \to q_4 \to q_5$ reconhece \texttt{\#} seguido de $\{\texttt{i},\texttt{s}\}^+$ (gerado por $T$ e $R$ via $P_5$--$P_9$), $q_5 \to q_6 \to q_7$ reconhece \texttt{:} seguido de $\{\texttt{i},\texttt{s}\}^+$ e $q_7 \to q_8$ consome \texttt{;}. Como $q_8$ não possui transições de saída, exige-se consumo integral da entrada. Logo $L(M_{\text{AP}}) = L(G)$.
